\documentclass[10pt]{scrartcl}
\usepackage[mdthm,colorsec,noauthor,noasy]{evan}
\usepackage{handout}
\DeclareMathOperator{\im}{range}
\title{MATH 110 --- Week 5}
\subtitle{Homework 5 --- Solutions}
\author{}\date{}\ihead{}
\begin{document}\maketitle
\begin{prob}\textbf{HW5 P1(a).}\par\idea{The two evaluations force factors $x$ and $x-1$; linear interpolation gives every output.}\begin{steps}\ii $Tp=0$ iff $p(0)=p(1)=0$. By the factor argument proved in Week 4, $p=x(x-1)q$ with $q\in\mathcal P_1(\mathbb R)$. Conversely every such polynomial vanishes at both points.\ii Thus $\ker T=\operatorname{span}(x(x-1),x^2(x-1))$. The two polynomials are independent by their different degrees, so they form a basis.\ii For arbitrary $(a,b)\in\mathbb R^2$, the polynomial $p=a+(b-a)x$ has $Tp=(a,b)$. Hence $T$ is onto.\end{steps}\end{prob}
\begin{prob}\textbf{HW5 P1(b,c).}\par\idea{The induced map records the two evaluations; cosets of $1,x$ already produce an output basis.}\begin{steps}\ii Let $K=\ker T$. The induced isomorphism is $\overline T(p+K)=(p(0),p(1))$. Its inverse sends $(a,b)$ to $a+(b-a)x+K$, by part (a).\ii The cosets $1+K,x+K$ map to $(1,1),(0,1)$, an independent spanning list of $\mathbb R^2$. An isomorphism preserves bases, so these two cosets are a quotient basis.\ii Equivalently, every $p+K$ equals $p(0)(1+K)+(p(1)-p(0))(x+K)$, and evaluations force these coefficients uniquely.\end{steps}\end{prob}
\begin{prob}\textbf{HW5 P2.}\par\idea{The containment $U\subseteq\ker T$ ensures changing the representative does not change its image under $T$.}\begin{steps}\ii If $v+U=w+U$, then $v-w\in U\subseteq\ker T$, so $Tv=Tw$. Thus $\widetilde T(v+U)=Tv$ is well-defined.\ii $\widetilde T(a(v+U)+(w+U))=T(av+w)=aTv+Tw$, proving linearity. Every quotient element is some $v+U$, so the prescribed formula also proves uniqueness.\ii A coset $v+U$ lies in the kernel iff $Tv=0$, or $v\in\ker T$. Thus $\ker\widetilde T=\{v+U:v\in\ker T\}=(\ker T)/U$.\end{steps}\end{prob}
\begin{prob}\textbf{HW5 P3.}\par\idea{A single nonzero output can be scaled to every field element.}\begin{steps}\ii If $\varphi$ is not zero, choose $v$ with $a=\varphi(v)\ne0$.\ii For each $b\in F$, $\varphi((b/a)v)=(b/a)\varphi(v)=b$. Hence $\varphi$ is surjective.\end{steps}\end{prob}
\begin{prob}\textbf{HW5 P4.}\par\idea{Dual basis values isolate one coefficient at a time.}\begin{steps}\ii Write $v=\sum_j a_jv_j$. Applying $\varphi_k$ gives $\varphi_k(v)=\sum_j a_j\delta_{kj}=a_k$, proving part (a).\ii For part (b), evaluate $\sum_j\psi(v_j)\varphi_j$ at $v_k$. The result is $\sum_j\psi(v_j)\delta_{jk}=\psi(v_k)$.\ii The two linear functionals agree on a basis, hence everywhere by Week 3 Proposition 21.\end{steps}\end{prob}
\begin{prob}\textbf{HW5 P5.}\par\idea{Entries of the original matrix are evaluations by the codomain dual basis.}\begin{steps}\ii Let $B=(v_1,\ldots,v_n)$ and $C=(w_1,\ldots,w_m)$, with dual bases $\varphi_j$ and $\psi_i$. Since $Tv_j=\sum_i a_{ij}w_i$, we have $(T^{\prime}\psi_i)(v_j)=\psi_i(Tv_j)=a_{ij}$.\ii The functional expansion gives $T^{\prime}\psi_i=\sum_j a_{ij}\varphi_j$, so entry $(j,i)$ of its matrix is $a_{ij}$.\ii This is exactly the transpose matrix $A^{\mathsf T}$.\end{steps}\end{prob}
\begin{prob}\textbf{HW5 P6.}\par\idea{Write the vanishing conditions on the two generators.}\begin{steps}\ii A row $(a,b,c)$ lies in $U^0$ iff $a+b=0$ and $b+c=0$, so it equals $b(-1,1,-1)$. Thus $U^0=\operatorname{span}((-1,1,-1))$ over $\mathbb R$.\ii The generators of $U$ are independent because the first coordinate of a zero relation forces the first coefficient zero, and the third coordinate forces the second zero. Hence $\dim U=2$.\ii The annihilator generator is nonzero, so $\dim U^0=1$. Therefore $\dim U+\dim U^0=2+1=3$.\end{steps}\end{prob}
\begin{prob}\textbf{HW5 P7 (optional).}\par\idea{Separate a nonzero vector by one functional, then use finite-dimensional equality of dimensions.}\begin{steps}\ii For fixed $v$, evaluation $j_V(v)$ is linear in the functional; the identity $j_V(av+w)(\varphi)=a\varphi(v)+\varphi(w)$ also makes $j_V$ linear.\ii If $v\ne0$, extend $v$ to a basis and define $\varphi$ to be $1$ on $v$ and zero on every other basis vector. Finite basis expansions define a linear functional; arbitrary-dimensional basis extension uses the cited choice result.\ii Then $j_V(v)(\varphi)=1$, so $j_V(v)\ne0$, proving zero kernel and injectivity.\ii If $V$ is finite-dimensional, $\dim V^{\prime\prime}=\dim V^{\prime}=\dim V$. The equal-dimensional injectivity criterion gives an isomorphism.\end{steps}\end{prob}
\begin{prob}\textbf{Exam practice, Question 2(a,b).}\par\idea{The dual map precomposes a functional with the original map.}\begin{steps}\ii Define $T^*(\varphi)=\varphi\circ T$ for $\varphi\in W^*$. It lies in $V^*$ since composition of linear maps is linear.\ii For $a\in\mathbb R$ and $\varphi,\psi\in W^*$, $T^*(a\varphi+\psi)(v)=a\varphi(Tv)+\psi(Tv)=(aT^*\varphi+T^*\psi)(v)$ for all $v$. Thus $T^*$ is linear.\ii For $\eta\in Z^*$, $(S\circ T)^*\eta=\eta\circ S\circ T=T^*(S^*\eta)$ by associativity, proving the product identity.\end{steps}\end{prob}
\begin{prob}\textbf{Exam practice, Question 2(c).}\par\idea{Surjectivity makes vanishing on all $Tv$ the same as vanishing on all of $W$.}\begin{steps}\ii If $T^*\varphi=0$, then $\varphi(Tv)=0$ for all $v\in V$.\ii For each $w\in W$, choose $v$ with $Tv=w$. Then $\varphi(w)=0$, so $\varphi$ is the zero functional.\ii Thus $\ker T^*=\{0\}$, and the zero-kernel criterion proves injectivity.\end{steps}\end{prob}\begin{prob}\textbf{Exam practice, Question 2(d).}\par\idea{Compose the given scalar functional with the map.}\begin{steps}\ii $(T^*\ell)(x,y)=\ell(T(x,y))=3(2x-y)=6x-3y$.\ii Thus $T^*\ell$ is the functional represented by row $(6,-3)$.\end{steps}\end{prob}
\end{document}
